% !TEX root = ../Rom-abcd.tex
%% %%%%%%%%%%%%%%%%%%%%%%%%%%%%%%%%%%%%%%%%%%%%%%%%%%%%%%%%%%%%%%%%%%%%
%% -- Testdatei für alle wesentlichen Features der Vorlage
%% -- content-ulgr.tex
%% -- Stand: 2026-03-07
%% %%%%%%%%%%%%%%%%%%%%%%%%%%%%%%%%%%%%%%%%%%%%%%%%%%%%%%%%%%%%%%%%%%%%
\renewcommand{\LongTitel}{Halbgruppen auf Banachräumen -- Ein Überblick}
\renewcommand{\ShortTitel}{Halbgruppen auf Banachräumen}
\renewcommand{\AutorenBeitrag}{Ulrich Groh}
%% -- Kapitelüberschrift und TOC-Eintrag
\addchap[\ShortTitel]{\LongTitel}
\addtocontents{toc}{\textsc{\AutorenBeitrag}}
\addtocontents{toc}{}
%% -- Kopfzeile
\markleft{\textsc{\AutorenBeitrag}}
\markright{\textsc{\ShortTitel}}
%% -- Autor
\begin{center}
    \textsc{\Large \AutorenBeitrag}\\[1.5em]
%    \includegraphics[width=8cm, height=8cm, keepaspectratio=true]
%        {./content-ulgr/autor-abcd}
\end{center}
%% -- Motto
\renewcommand{\dictumwidth}{0.45\textwidth}
\dictum[\href{https://de.wikipedia.org/wiki/Israel_Gelfand}{Israel Gelfand}]{%
    Alle Mathematik ist Funktionalanalysis.}
%% -- Zusammenfassung
\begin{quote}
Dieser Beitrag dient als Testdatei für alle wesentlichen Features der
Vorlage: Theorem-Umgebungen, mathematische Notation, Abkürzungen,
Literaturverweise, Aufzählungen und Tabellen.
\end{quote}
%% %%%%%%%%%%%%%%%%%%%%%%%%%%%%%%%%%%%%%%%%%%%%%%%%%%%%%%%%%%%%%%%%%%%%
\begin{refsection}
%% %%%%%%%%%%%%%%%%%%%%%%%%%%%%%%%%%%%%%%%%%%%%%%%%%%%%%%%%%%%%%%%%%%%%

%% %%%%%%%%%%%%%%%%%%%%%%%%%%%%%%%%%%%%%%%%%%%%%%%%%%%%%%%%%%%%%%%%%%%%
\section*{Test der Vorlage}
%% %%%%%%%%%%%%%%%%%%%%%%%%%%%%%%%%%%%%%%%%%%%%%%%%%%%%%%%%%%%%%%%%%%%%

\subsection*{Abkürzungen}

Einige Abkürzungen im Fließtext: \zB ist $E = L^2([0,1])$ ein
Hilbertraum; \bzw ist $E = C([0,1])$ ein Banachraum.
Die Theorie gilt \zB für beliebige Banachräume, \vgl \textcite{alsina:2013}.
Weitere Details findet man \ua in \textcite{takesaki:1958} \bzw in
den dort angegebenen Quellen.

\subsection*{Mathematische Symbole}

Zahlbereiche: $\N, \Z, \Q, \R, \C, \K$.

Norm und Betrag: $\norm{x}$, $\norm{}$, $\abs{x}$, $\abs{}$.

Differentiale: $\int_0^t f(s)\ds$, $\int_\R f(x)\dx$,
$\int_0^\infty \eu^{-\lambda t} T(t)x\dt$.

Eulersche Zahl und imaginäre Einheit: $\eu^{\iu\pi} + 1 = 0$.

\subsection*{Theorem-Umgebungen ohne Nummer}

\begin{theorem}[Ohne Nummer]
Dies ist ein Theorem ohne Nummer -- \verb|\begin{theorem}|.
\end{theorem}

\begin{lemma}
Dies ist ein Lemma ohne Nummer -- \verb|\begin{lemma}|.
\end{lemma}

\begin{remark}
Dies ist eine Anmerkung -- \verb|\begin{remark}|.
In \texttt{\small} gesetzt.
\end{remark}

\begin{example}
Auf $E = C_0(\R)$ wirkt die Translationshalbgruppe durch
$T(t)f(x) = f(x+t)$ -- \verb|\begin{example}|.
\end{example}

\subsection*{Theorem-Umgebungen mit Nummer}

\begin{ntheorem}[Mit Nummer]
Dies ist ein Theorem ohne Nummer -- \verb|\begin{theorem}|.
\end{ntheorem}

\begin{nlemma}
Dies ist ein Lemma ohne Nummer -- \verb|\begin{lemma}|.
\end{nlemma}

\begin{nremark}
Dies ist eine Anmerkung -- \verb|\begin{remark}|.
In \texttt{\small} gesetzt.
\end{nremark}

\begin{nexample}
Auf $E = C_0(\R)$ wirkt die Translationshalbgruppe durch
$T(t)f(x) = f(x+t)$ -- \verb|\begin{example}|.
\end{nexample}

\subsection*{Aufzählungen}

Eine Aufzählung mit \texttt{enumerate}:
\begin{enumerate}[\upshape(i)]
    \item Erster Punkt,
    \item zweiter Punkt,
    \item dritter Punkt.
\end{enumerate}

Eine Aufzählung mit \texttt{itemize}:
\begin{itemize}
    \item $C_0$-Halbgruppen,
    \item Kontraktionshalbgruppen,
    \item analytische Halbgruppen.
\end{itemize}

\subsection*{Tabelle}

\begin{center}
\begin{tabular}{lll}
    \toprule
    Raum $E$ & Halbgruppe $T(t)$ & Erzeuger $A$ \\
    \midrule
    $L^{p}(\R)$ & Translation $f(\,\cdot\,+t)$ & Ableitung $f'$ \\
    $L^{2}([0,1])$ & Wärme & $\Delta$ mit Randbedingungen \\
    $C_{0}(\R)$ & Translation & Ableitung $f'$ \\
    \bottomrule
\end{tabular}
\end{center}

\subsection*{Anführungszeichen}

Deutsch: \enquote{Ein schöner Satz.} \verb|\enquote{Ein schöner Satz.}|

Englisch: \foreignquote{english}{A beautiful theorem.} \verb|\foreignquote{english}{A beautiful theorem.}|

%% %%%%%%%%%%%%%%%%%%%%%%%%%%%%%%%%%%%%%%%%%%%%%%%%%%%%%%%%%%%%%%%%%%%%
%% -- Literaturverzeichnis
%% %%%%%%%%%%%%%%%%%%%%%%%%%%%%%%%%%%%%%%%%%%%%%%%%%%%%%%%%%%%%%%%%%%%%
\RaggedRight
\printbibliography
\end{refsection}
%% -- Ende des Beitrags
